\chapter{Dust grains and grain populations}\label{ch:dust}
\section{Size is a physical state}
Dust is not a gas species with a larger fixed mass. Radius, material density,
charge, velocity and any optical/material properties control its dynamics.
For an ideal sphere of radius $a$ and density $\rho_d$,
\begin{equation}
m_d=\frac{4\pi}{3}\rho_da^3,\qquad
\phi=\frac{q_d}{4\pi\epsilon_0a},\qquad
\boldsymbol a_L=\frac{q_d}{m_d}
(\boldsymbol E+\boldsymbol v\times\boldsymbol B).
\end{equation}
The potential relation assumes an isolated spherical conductor. Shielding,
nonspherical grains and material effects require different closures.
The Lorentz term includes neither gas drag, gravity nor radiation pressure.

\section{Size distributions and importance weights}
For a number distribution $d\dot N/da=Ca^{-p}$ on
$[a_{\min},a_{\max}]$, total number and mass rates are different integrals:
\begin{align}
\dot N&=C\int a^{-p}\,da,\\
\dot M&=\frac{4\pi\rho_dC}{3}\int a^{3-p}\,da .
\end{align}
The logarithmic limits $p=1$ and $p=4$ must be evaluated with their proper
logarithms. A specified dust mass source does not determine a numerical
particle weight until the size proposal and material density are given.
For proposal density $g(a)$, an event's represented contribution contains
the physical density divided by $g(a)$. This correction is independent of
generic population splitting/merging.

The PIC-coupled dust model requires positive individual weight corrections.
Logarithmic radius bins are half-open, with the global upper endpoint owned
by the final bin. A steep number distribution can coexist with most of the
injected mass in the largest grains. The in-coma size distribution additionally
depends on size-dependent speed and survival; it need not reproduce the
source distribution.

\section{Charging and forces}
Charging integrates currents appropriate to the declared plasma/illumination
model. Validate equilibrium roots and transient convergence across multiple
radii, including the sign convention for collection and emission currents.
A floating potential can be small while charge-to-mass acceleration is large.
If plasma inputs are absent, an explicit neutral/uncharged model is required;
zero-filled background data are not a physically justified charging solution.

Gas drag, radiation pressure, gravity and Lorentz forces have distinct
validity regimes. Couple each force to the same grain state and frame.
Radiation/optical interpolation tables need wavelength, radius, refractive-
index convention and out-of-range policy. A Mie table lookup test is not
validation of an entire dust tail.

\section{Dust-specific split and merge}
The dependency-light dust API groups particles by log radius and requires the
same native owner cell, compiled species and size group. Current merging
additionally requires exactly equal intrinsic radius, mass and charge.
Same-bin membership alone is insufficient: two radii can have different
charging and drag physics even if their difference looks small numerically.

Splitting clones one grain into $n$ representatives with weights $w/n$.
The supported merge maps at least three compatible particles to two:
total weight $W$, weighted mean velocity $\boldsymbol u$, and velocities
$\boldsymbol u\pm d\boldsymbol e_j$, each with weight $W/2$, where
$d^2$ is the weighted velocity variance. It conserves represented number,
mass, charge, momentum and scalar kinetic energy. It does not generally
preserve the pressure tensor, higher moments or mean speed. The API can
reject a requested sampled-speed constraint transactionally.

\scope{The dust split/merge physical kernel is implemented and tested.
A general production PIC-list adapter is not supplied by that algebraic
PASS. Generic gas split/merge must not be used to erase radius-dependent
physics.}

\section{Spherical injection and its ledger}
The production sphere injector chooses surface position relative to the
actual centre, applies native ownership and local-time acceptance, creates
the record and dispatches the selected mover. A signed residual mass reservoir
accounts for finite particle granularity. A negative residual is not negative
physical grain mass; it records that an event has advanced a fractional budget.
The native fixture checks this ledger over two intervals and verifies that
gas NTC leaves steep-weight grain records unchanged.
\source{src/models/dust/README.md; test/srcModelDustTest/README.md; test/srcModelDustInjectionTest/README.md}
